\chapter[Stellung der Synode Iowa zu den Symbolen]{„Stellung der Synode Iowa zu den Symbolen der evangelisch-lutherischen Kirche“}
\quellseite{29}

Unter dieser Ueberschrift legt die Gemeinschaft derjenigen, welche nicht nur durch Vermittlung Herrn Pfarrer Löhe’s nach Amerika gekommen sind,
\quellseite{30}
sondern demselben auch unter allen Umständen ergeben bleiben wollen, ihre Stellung zum Concordienbuch dar. Es hat diese Darstellung namentlich darum Interesse, weil sie zeigt, wie verschieden das Thema unserer Zeit, daß man die kirchlichen Bekenntnisse nicht mehr unqualificirt annehmen könne, sich ausdrücken lasse; ferner weil diese Jowa-Synode eine Schutzgenossin der Buffalo-Synode; und endlich weil sie eben Herrn Pfarrer Löhe’s folgsames Töchterchen ist oder wenigstens von demselben dafür angesehen und deswegen reichlich mit Geldmitteln versehen wird. Vorläufig können wir nur unser tiefes Bedauern darüber aussprechen, daß in dem Document nicht Herrn Pfarrer Löhe’s Ehrlichkeit und Aufrichtigkeit zu sehen ist, da die Schreiber auf der einen Zeile versichern, sich „rückhaltslos“ zur Concordia von 1580 zu bekennen, ja sie als „norma docendi“, Regel und Richtschnur der Lehre anzusehen, auf der anderen Seite vor einer „vorwiegend dogmatischen Auffassung“ warnen, sich für eine f. g. „historische Auffassung“ erklären und namentlich manche Beweisführungen und Erklärungen, die darin vorkommen, als zum Bekenntniß nicht gehörig von sich weisen (z.\,B. notorisch dem Chiliasmus huldigen und die Entwicklung der Lehre von Kirche und Amt mit Pfarrer Löhe für unrichtig halten). Doch davon, f. G. w., später, da wir das Jowaer Kirchenblatt eben erst in erster und zweiter Nummer unmittelbar vor Schluß dieses Heftes erhielten. Die Worte unter obigem Titel lauten wie folgt:

1. Wir erkennen die Concordia von 1580 in allen ihren Theilen für unser Bekenntniß, auf welchem unsere Pastoren und Gemeinden stehen und, weil sie der luth. Kirche angehören, stehen müssen.

2. Alle unsere Pastoren werden deshalb bei der Ordination und Installation ohne Rückhalt auf die Concordia von 1580 verpflichtet. Eben so bekennen sich alle unsere Gemeinden, laut unserer Kirchenordnung, rückhaltslos zu derselben. Bekenntnißuntreue Gemeinden werden von uns nicht pastorirt, noch in unsere Synodalgemeinschaft aufgenommen.

3. Wir erkennen die Symbole an als das, was sie sind und laut der Concordienformel selber sein wollen: als Zeugnisse der rechten und reinen Erkenntniß und Lehre des Wortes Gottes, welche Gott seiner Kirche in den mannigfaltigen vor und in der Reformationszeit aufgekommenen Streitigkeiten gegeben, welche sie als ihre anvertraute, seligmachende Wahrheit festgehalten und gegen Lüge und Irrthum vertheidiget hat. Eben deshalb bilden sie die norma docendi, Regel und Richtschnur der Lehre.

4. Weil die Symbole größtentheils Resultate kirchlichen Kampfes sind und Entscheidung und Beilegung des Kampfes zum Zweck und Ziel haben, so können wir allein die historische Auffassung für die richtige, der Natur der Bekenntnisse entsprechende erkennen. Denn nur aus dem jedesmaligen Streit kann eine Bekenntnißschrift und ein Artikel derselben richtig verstanden und aufgefaßt werden. Jene vorwiegend dogmatische, unhistorische Auffassung, welche die historische übersteht und verkennt, durch welche die aufgestellten Bekenntnißsätze vielfach erst ins Licht treten und ihre richtige Begrenzung erhalten, müssen wir für eine unrichtige und verkehrte erklären.

5. Da unsere Bekenntnißschriften zum Theil nicht bloß die einfachen Resultate des kirchlichen Kampfes, d.\,h. nicht bloß die bekennenden Sätze und verwerfenden Gegensätze enthalten, sondern dieselben auch erläutern, begründen und vertheidigen: so ist allerdings ein Unterschied zu machen zwischen Wesentlichem und Unwesentlichem in diesen Schriften. (Eigentliches Bekenntniß, die Gewissen bindende norma docendi können bloß die thetischen
\quellseite{31}
und antithetischen Entscheidungen sein, welche jeder Artikel der Lüge und dem Irrthum gegenüber ausspricht und feststellt. Dagegen kann nicht jede Beweisführung, jede Erläuterung 2c., die eigentlich in einer Bekenntnißschrift Rechtens ist, zu einem die Gewissen bindenden Glaubenssatz gemacht werden. Symbolische Geltung hat, was die Symbole symbolisch feststellen wollen, und das liegt in jedem Artikel für den, der ihn im Lichte der Historie liest, auf platter Hand.

6. Wir glauben allerdings nicht, daß mit der Reformationszeit die Lehrentwicklung der Kirche zu ihrem Ende gekommen sei; wir glauben, daß noch Lehrkämpfe aufkommen können und vielleicht auch werden, die ihre Entscheidung in der Reformationszeit noch nicht finden konnten. Darauf gründet sich der Satz in unserer Kirchenordnung, daß wir uns „zu derjenigen Richtung der luth. Kirche bekennen, welche auf dem Wege der Symbole an der Hand des Wortes Gottes einer größeren Vollendung der evang.-luth. Kirche entgegenstrebt.“

Wir sind der festen Ueberzeugung, daß auch die Lehre von der Kirche und dem Amte, so weit sie in Symbolen (dem römischen Gegensatz gegenüber) ihre Bestimmungen gefunden hat, dem Worte Gottes entsprechend ist, und erkennen die Symbole auch in diesem Stück, wie in allen andern, von welchen sie handeln, für unsere norma docendi.
